\section{Discussion}

\paragraph{Industrial Deployment and Operational Impact.}
Our approach enables selective deployment decisions at the commit level during live events, potentially reducing freeze scope and accelerating feature velocity. Deployment freezes currently block 100\% of changes regardless of risk, though only approximately 5\% of commits historically cause incidents. On held-out 2024 production data, the model achieved 0.926 precision and 0.788 recall, suggesting it could clear the vast majority of safe changes for deployment while flagging most risky ones for human review---substantially reducing freeze scope without compromising safety. The threshold is tunable: lowering it increases recall further at modest precision cost, acceptable for a system that flags for review rather than auto-blocks. Adoption barriers include CI/CD integration, threshold calibration against business impact, and developer trust through transparent explanations. Validating these estimates through shadow-mode deployment during live events is planned as immediate next work.

\paragraph{Trade-offs and Design Constraints.}
While combining diff-aware features with historical metadata could enhance performance, Prime Video's operational constraints make a diff-only approach optimal: microservice churn makes historical data stale, cross-functional onboarding creates cold-start scenarios, and diff-aware features preserve privacy. Our ablation study (F1: 0.846) suggests diminishing returns from metadata enrichment may not justify maintenance costs.

A key practical benefit of LLM-based feature extraction is eliminating language-specific static analysis tooling. Traditional approaches require separate AST parsers and extraction scripts for each language, each maintained independently as languages evolve. A single LLM extractor generalizes across languages with one prompt template, significantly lowering the cost of onboarding new languages or repositories.

\paragraph{Feature Contribution Insights.}
A counterintuitive finding emerged from our ablation study: quantitative change-level features (lines added/deleted, chunks modified) degraded model performance, suggesting raw change volume is a noisy proxy for risk. Code-level structural complexity and qualitative indicators provide more reliable signals, confirmed by the ablation results (complexity-only F1: 0.809 vs.\ size-only F1: 0.565).

These LLM baselines represent the state of the art for end-to-end agentic risk classification without feature engineering. Performance improved from F1 0.153 to 0.231 (zero-shot) and 0.216 to 0.261 (few-shot) across model generations, demonstrating that newer models reduce the gap but do not close it---structured feature curation remains essential. This gap underscores that decomposing diffs into interpretable quantitative and qualitative dimensions captures risk patterns that end-to-end LLM classification misses, while also producing explainable predictions suitable for developer-facing guardrails.

ApacheJIT's lower F1-score (0.771 vs 0.846) likely stems from dataset characteristics: Prime Video's CoE-based labeling versus ApacheJIT's diverse post-release defect labeling. Our findings extend prior observations that code-level features carry substantial predictive power~\cite{kamei2012large,hoang2019deepjit} to a diff-only standalone model for cold-start scenarios~\cite{fukushima2014cross,liu2025human}.
